\documentclass[12pt]{article}
\usepackage{latexblog}

% ============================================================
% Blog Metadata
% ============================================================
\blogtitle{IEEE 754浮点数转换}
\blogdate{2019-06-27}
\blogtags{刨根问底, 协议}
\bloglang{zh}
\blogsummary{}

\begin{document}
\makeblogtitle

一个小数的二进制是怎么样的呢？我们先看看一个二进制的小数怎么转换成十进制： \[
\begin{aligned} 
11101.01011_{10} &= 1 \times 2^{4} + 1 \times 2^{3} + 1 \times 2^{2} + 0 \times 2^{1} + 1 \times 2^{0} + 0 \times 2^{-1} + 1 \times 2^{-2} + 1 \times 2^{-3} + 1 \times 2^{-4} + 1 \times 2^{-5} \\
  &= 16 + 8 + 4 + 0 + 1 + 0 + \frac{1}{2} + 0 + \frac{1}{16} + \frac{1}{32} \\
  &= 29.34375
\end{aligned}
\]

\section{IEEE 754}\label{ieee-754}

\href{https://en.wikipedia.org/wiki/IEEE_754}{IEEE 754} 标准中规定了浮点数在计算机中的表示方法，主要就是单精度(float)和双精度(double):

\begin{verbatim}
              S Exp      Fraction
Single(32bit) ▯▮▮▮▮▮▮▮▮▯▯▯▯▯▯▯▯▯▯▯▯▯▯▯▯▯▯▯▯▯▯▯
Double(64bit) ▯▮▮▮▮▮▮▮▮▮▮▮▯▯▯▯▯▯▯▯▯▯▯▯▯▯▯▯▯▯▯▯▯▯▯▯▯▯▯▯▯▯▯▯▯▯▯▯▯▯▯▯▯▯▯▯▯▯▯▯▯▯▯▯
\end{verbatim}

其计算公式为： \[
x = (-1)^{S}\times(1+ Fraction)\times 2^{(Exponent-Bias)}
\] 其中，Bias为\href{https://zh.wikipedia.org/wiki/IEEE_754\#\%E6\%8C\%87\%E6\%95\%B8\%E5\%81\%8F\%E7\%A7\%BB\%E5\%80\%BC}{指数偏移值}，是一个固定值，即\(Bias=2^{e-1} - 1\) 其中e为指数部分的比特长度。

\begin{itemize}
\tightlist
\item
  单精度\(Bias = 2^{7} - 1 = 127\)
\item
  双精度\(Bias = 2^{10} -1 = 1023\)
\end{itemize}

举个例子，刚才我们算出来的小数可以这样表示：

\[
\begin{aligned} 
29.34375 &= 11101.01011 = 11101.01011_{2} \times 2^{0} \\
    &= 1.110101011_{2} \times 2^{4} \\
    &= (-1)^{0} \times (1 + 0.110101011_{2}) \times 2^{131 - 127} \\
    &= (-1)^{0} \times (1 + 0.110101011_{2}) \times 2^{10000011_{2} - 127} \\
    &= (-1)^{0} \times (1 + 0.110101011_{2}) \times 2^{1027 - 1023} \\
    &= (-1)^{0} \times (1 + 0.110101011_{2}) \times 2^{10000000011_{2} - 127} \\
\end{aligned}
\]

因此在计算机中，表示为：

\begin{verbatim}
Float:
▯▮▮▮▮▮▮▮▮▯▯▯▯▯▯▯▯▯▯▯▯▯▯▯▯▯▯▯▯▯▯▯
010000011110101011..............

Double:
▯▮▮▮▮▮▮▮▮▮▮▮▯▯▯▯▯▯▯▯▯▯▯▯▯▯▯▯▯▯▯▯▯▯▯▯▯▯▯▯▯▯▯▯▯▯▯▯▯▯▯▯▯▯▯▯▯▯▯▯▯▯▯▯
010000000011110101011...........................................
\end{verbatim}

不足的部分补上0，即：\(0100000111101010110000000000000_{2}=41eac000{16}\)

除了正常的浮点数外，还有几个比较特殊的：

{\def\LTcaptype{none} % do not increment counter
\begin{longtable}[]{@{}ll@{}}
\toprule\noalign{}
Description & Float(32bit) \\
\midrule\noalign{}
\endhead
\bottomrule\noalign{}
\endlastfoot
Zero & 0 00000000 00000000000000000000000 \\
Negative Zero & 1 00000000 00000000000000000000000 \\
Infinity & 0 11111111 00000000000000000000000 \\
Negative Infinity & 1 11111111 00000000000000000000000 \\
Not a Number (NaN) & 0 11111111 00001000000000100001000 \\
\end{longtable}
}

\section{十进制与二进制转换}\label{ux5341ux8fdbux5236ux4e0eux4e8cux8fdbux5236ux8f6cux6362}

计算方式为，将小数的整数部分与2取余倒序排列；将小数部分与2取整正序排列。例如，将3.14转换为float:

\begin{itemize}
\tightlist
\item
  首先将整数部分直接转换为二进制 \(3_{10} = 11_{2}\)

  \begin{itemize}
  \tightlist
  \item
    \(3\mod2 = \fbox{1}\)
  \item
    \(1\mod2 = \fbox{1}\)
  \end{itemize}
\item
  小数部分为0.14，不断乘以2后取整数部分，然后用小数继续乘以2直到值为1

  \begin{itemize}
  \tightlist
  \item
    \(0.14 \times 2 = \fbox{0}.28\)
  \item
    \(0.28 \times 2 = \fbox{0}.56\)
  \item
    \(0.56 \times 2 = \fbox{1}.12\)
  \item
    \(0.12 \times 2 = \fbox{0}.24\)
  \item
    \(0.24 \times 2 = \fbox{0}.48\)
  \item
    \(0.48 \times 2 = \fbox{0}.96\)
  \item
    \(0.96 \times 2 = \fbox{1}.92\)
  \item
    \(0.92 \times 2 = \fbox{1}.84\)
  \item
    \ldots..
  \item
    重复以上步骤，得到\(0.14_{10}=0.001000111101011100001010001111010..._{2}\)
  \end{itemize}
\end{itemize}

\section{舍入操作}\label{ux820dux5165ux64cdux4f5c}

对于尾数多余精度的情况，需要舍去多余的部分，但不是按照四舍五入的方式，而是按照''向偶舍入``的方式，意思就是，如果多余的部分大于0.5(\(0.5_{10} = 0.1_{2})\)则最低位进1；如果小于0.5则舍去；如果正好是等于0.5则根据最低位判断，如果最低位是1则进位，否则舍去。这样按照统计学来看，对于一个小数有相同的机会进位或者被舍去。

例如对于上例中的3.14，我们可以得到：

\[
\begin{aligned} 
3.14 &= 11.001000111101011100001010001111010..._{2} \\
    &= 1.1001000111101011100001010001111010..._{2} \times 2^{1} \\
    &= (-1)^{0} + (1 + 0.1001000111101011100001010001111010..._{2}) \times 2^{128 - 127} \\
    &= (-1)^{0} + (1 + 0.10010001111010111000010{\color{blue}{10001111010...}}_{2}) \times 2^{10000000_{2} - 127} \\
    &\approx (-1)^{0} + (1 + 0.1001000111101011100001{\color{red}1}_{2}) \times 2^{10000000_{2} - 127}
\end{aligned}
\]

因此3.14的float表示为：

\begin{verbatim}
▯▮▮▮▮▮▮▮▮▯▯▯▯▯▯▯▯▯▯▯▯▯▯▯▯▯▯▯▯▯▯▯
01000000010010001111010111000011
\end{verbatim}

\section{还原}\label{ux8fd8ux539f}

那么，对于一个存储在磁盘上的浮点数，我们怎么将它加载到内存中来？对于C来说，实际也是采用的IEEE754标准(float, double)，所以实际上浮点数在内存中的表示是一致的，直接转换即可：

\begin{Shaded}
\begin{Highlighting}[]
\KeywordTok{struct}\NormalTok{ ConstantFloat }\OperatorTok{\{}
        \AttributeTok{mutable}\NormalTok{ u4 bytes}\OperatorTok{;}
        \DataTypeTok{float} \OperatorTok{\&}\NormalTok{getValue}\OperatorTok{()} \AttributeTok{const} \OperatorTok{\{}
            \ControlFlowTok{return} \OperatorTok{*}\KeywordTok{reinterpret\_cast}\OperatorTok{\textless{}}\DataTypeTok{float} \OperatorTok{*\textgreater{}(\&}\NormalTok{bytes}\OperatorTok{);}
        \OperatorTok{\}}
    \OperatorTok{\};}
\end{Highlighting}
\end{Shaded}

参考：

\begin{itemize}
\tightlist
\item
  \href{http://sandbox.mc.edu/~bennet/cs110/flt/dtof.html}{Decimal to Floating-Point Conversions}
\item
  \href{http://cs.boisestate.edu/~alark/cs354/lectures/ieee754.pdf}{IEEE 754 FLOATING POINT REPRESENTATION}
\item
  \href{https://www.h-schmidt.net/FloatConverter/IEEE754.html}{IEEE-754 Floating Point Converter}
\item
  \href{https://www.rapidtables.com/convert/number/binary-to-decimal.html}{Binary to Decimal converter}
\item
  \href{https://www.jianshu.com/p/e5d72d764f2f}{IEEE754表示浮点数}
\item
  \href{http://www.binaryconvert.com/result_double.html?decimal=050057046051052051055053}{Online Binary-Decimal Converter}
\end{itemize}


\end{document}
